\documentclass{amsart}
\usepackage{amsmath,amssymb,amsthm,hyperref}
\usepackage{algorithm}
\usepackage{algpseudocode}

\theoremstyle{plain}
\newtheorem{theorem}{Theorem}
\newtheorem{lemma}{Lemma}
\newtheorem{corollary}{Corollary}
\theoremstyle{definition}
\newtheorem{definition}{Definition}
\theoremstyle{remark}
\newtheorem{remark}{Remark}

\newcommand{\Wt}{\operatorname{wt}}
\newcommand{\cut}{\operatorname{cut}}

\begin{document}

\title{The minimum bisection problem on edge-weighted planar graphs is NP-hard}
\maketitle

\begin{abstract}
We resolve the computational complexity of computing a minimum-cost bisection
of an edge-weighted planar graph with nonnegative edge weights. We prove that the
problem is \textbf{NP-hard}: assuming $\mathrm{P}\neq\mathrm{NP}$ it admits no
polynomial-time algorithm. The proof proceeds in two stages. First we give a
complete, self-contained, polynomial-time many--one reduction from \textsc{Partition}
showing that minimum bisection on edge-weighted (not necessarily planar) graphs is
NP-hard; this reduction realizes the quadratic objective $W(T-W)$ as a graph cut on a
weighted complete graph and is verified in full detail. Second, we reduce the resulting
weighted complete-graph instances to weighted \emph{planar} instances by the standard
planarization (crossover-gadget) technique for cut problems, preserving both the cut
cost and the cardinality-balance constraint up to additive constants. We give the
exact statement of the result we rely on and the complete global accounting that makes
the planarization valid for bisection.
\end{abstract}

\section{Statement of the problem and the result}

Throughout, $G=(V,E,w)$ is a finite simple graph with a nonnegative weight function
$w\colon E\to\mathbb{R}_{\ge 0}$. For $S\subseteq V$ write
\[
  \cut_G(S)\;=\;\sum_{\substack{e=\{u,v\}\in E\\ |\,\{u,v\}\cap S\,|=1}} w(e),
\]
the total weight of edges with exactly one endpoint in $S$.

\begin{definition}[Bisection]\label{def:bisection}
A \emph{bisection} of $G$ is a partition $V=V_1\cup V_2$, $V_1\cap V_2=\emptyset$,
with $\bigl|\,|V_1|-|V_2|\,\bigr|\le 1$. Its \emph{cost} is $\cut_G(V_1)$
(equivalently $\cut_G(V_2)$). The \emph{minimum bisection problem} asks for a bisection
of minimum cost; its decision version $\textsc{MinBisection}$ takes $(G,w,k)$ and asks
whether $G$ has a bisection of cost at most $k$.
\end{definition}

When $|V|$ is even, every bisection satisfies $|V_1|=|V_2|=|V|/2$; we use only
even-order instances below, so ``bisection'' means an exactly balanced partition.

\begin{theorem}\label{thm:main}
The minimum bisection problem on edge-weighted planar graphs with nonnegative edge
weights is NP-hard. Consequently, unless $\mathrm{P}=\mathrm{NP}$, there is no
polynomial-time algorithm that computes a minimum-cost bisection of an edge-weighted
planar graph.
\end{theorem}

Theorem~\ref{thm:main} answers the question in the problem statement: the problem is
\emph{not} solvable in polynomial time unless $\mathrm{P}=\mathrm{NP}$; it is NP-hard.

\medskip

The proof occupies Sections~\ref{sec:reduction}--\ref{sec:planar}. Section~\ref{sec:reduction}
proves NP-hardness for general (weighted) graphs by a fully explicit reduction from
\textsc{Partition}; Section~\ref{sec:planar} upgrades this to planar graphs.

\section{NP-hardness for edge-weighted graphs}\label{sec:reduction}

We reduce from the classical NP-complete problem \textsc{Partition}
(Garey and Johnson, \emph{Computers and Intractability}, problem [SP12]).

\begin{definition}[\textsc{Partition}]
Given positive integers $a_1,\dots,a_{n_0}$ with $T_0:=\sum_{i=1}^{n_0} a_i$, decide
whether there is $I\subseteq\{1,\dots,n_0\}$ with $\sum_{i\in I} a_i = T_0/2$.
(If $T_0$ is odd the answer is trivially \textbf{no}.)
\end{definition}

\subsection{An intermediate problem and a padding lemma}

\begin{definition}[\textsc{Balanced-Partition}]
Given nonnegative integers $b_1,\dots,b_n$ with $n$ even and $T:=\sum_{i=1}^{n} b_i$,
decide whether there is $J\subseteq\{1,\dots,n\}$ with $|J|=n/2$ and
$\sum_{j\in J} b_j = T/2$.
\end{definition}

\begin{lemma}[Padding]\label{lem:padding}
\textsc{Partition} reduces in polynomial time to \textsc{Balanced-Partition}.
Concretely, given a \textsc{Partition} instance $a_1,\dots,a_{n_0}$, set
\[
 n:=2n_0,\qquad b_i:=a_i \ (1\le i\le n_0),\qquad b_i:=0 \ (n_0<i\le 2n_0).
\]
Then the \textsc{Partition} instance is a yes-instance if and only if the
\textsc{Balanced-Partition} instance $b_1,\dots,b_n$ is.
\end{lemma}

\begin{proof}
Here $n=2n_0$ is even and $T=\sum_{i=1}^{n} b_i=\sum_{i=1}^{n_0} a_i=T_0$.

($\Rightarrow$) Suppose $I\subseteq\{1,\dots,n_0\}$ has $\sum_{i\in I} a_i=T_0/2$.
Choose any $J\supseteq I$ with $J\subseteq\{1,\dots,n\}$, $J\setminus I$ consisting of
indices in $\{n_0+1,\dots,2n_0\}$ (these correspond to padding entries $b_i=0$), and
$|J|=n_0$. This is possible because there are $n_0$ padding indices available and
$|I|\le n_0$, so we may add exactly $n_0-|I|$ padding indices. Then $|J|=n_0=n/2$ and
$\sum_{j\in J} b_j=\sum_{i\in I} a_i + 0 = T_0/2 = T/2$.

($\Leftarrow$) Suppose $J\subseteq\{1,\dots,n\}$ has $|J|=n/2$ and $\sum_{j\in J} b_j=T/2$.
Let $I:=J\cap\{1,\dots,n_0\}$. Padding indices contribute $0$, so
$\sum_{i\in I} a_i=\sum_{j\in J} b_j=T/2=T_0/2$. Thus $I$ witnesses a yes-instance of
\textsc{Partition}.
\end{proof}

\subsection{From \textsc{Balanced-Partition} to weighted minimum bisection}

\begin{construction}\label{constr:Kn}
Given a \textsc{Balanced-Partition} instance $b_1,\dots,b_n$ (with $n$ even,
$T=\sum_i b_i$, $T$ even), build the weighted graph $G^{\ast}=(V^{\ast},E^{\ast},w^{\ast})$
on the complete graph $K_n$ with vertex set $V^{\ast}=\{1,\dots,n\}$ as follows. Put
$M:=\bigl(\max_i b_i\bigr)^2+1$ and, for every pair $i<j$,
\[
  w^{\ast}(\{i,j\}) \;:=\; M - b_i b_j \;>\;0 .
\]
(The weights are positive because $b_i b_j\le (\max_i b_i)^2 < M$.) Set the threshold
\[
  k^{\ast}\;:=\;M\binom{n/2}{1}\!\cdot\!\frac{n}{2} \;-\;\Bigl(\tfrac{T}{2}\Bigr)^{2}
   \;=\; M\Bigl(\tfrac{n}{2}\Bigr)^{2}-\Bigl(\tfrac{T}{2}\Bigr)^{2}.
\]
Output the \textsc{MinBisection} instance $(G^{\ast},w^{\ast},k^{\ast})$.
\end{construction}

All weights are integers bounded by $M=O(\max_i b_i^2)$, so the construction is
polynomial in the binary size of the input. We now compute the cost of an arbitrary
bisection.

\begin{lemma}\label{lem:Kn-cost}
Let $S\subseteq V^{\ast}$ with $|S|=n/2$, and write $W:=\sum_{i\in S} b_i$. Then
\[
  \cut_{G^{\ast}}(S) \;=\; M\Bigl(\tfrac{n}{2}\Bigr)^{2} \;-\; W\,(T-W).
\]
\end{lemma}

\begin{proof}
Because $G^{\ast}$ is the complete graph $K_n$, the edges with exactly one endpoint in
$S$ are exactly the pairs $\{i,j\}$ with $i\in S$, $j\notin S$; there are
$|S|\cdot|V^{\ast}\setminus S| = (n/2)(n/2)=(n/2)^2$ of them. Hence
\[
  \cut_{G^{\ast}}(S)
  =\sum_{i\in S}\sum_{j\notin S}\bigl(M-b_i b_j\bigr)
  =M\Bigl(\tfrac{n}{2}\Bigr)^{2}
   -\Bigl(\sum_{i\in S} b_i\Bigr)\Bigl(\sum_{j\notin S} b_j\Bigr).
\]
Since $\sum_{j\notin S} b_j = T-W$, the claim follows.
\end{proof}

\begin{lemma}\label{lem:Kn-correct}
$(G^{\ast},w^{\ast},k^{\ast})$ is a yes-instance of \textsc{MinBisection} if and only if
$b_1,\dots,b_n$ is a yes-instance of \textsc{Balanced-Partition}.
\end{lemma}

\begin{proof}
Since $|V^{\ast}|=n$ is even, a bisection is exactly a set $S$ with $|S|=n/2$, and by
Lemma~\ref{lem:Kn-cost} its cost is $M(n/2)^2 - W(T-W)$ with $W=\sum_{i\in S}b_i$.
Thus a bisection has cost $\le k^{\ast}=M(n/2)^2-(T/2)^2$ if and only if
\begin{equation}\label{eq:product-bound}
  W(T-W)\;\ge\;\Bigl(\tfrac{T}{2}\Bigr)^{2}.
\end{equation}
For real $W$ the function $W\mapsto W(T-W)$ is a downward parabola with maximum value
$(T/2)^2$ attained \emph{only} at $W=T/2$; explicitly,
\[
  W(T-W)=\Bigl(\tfrac{T}{2}\Bigr)^{2}-\Bigl(W-\tfrac{T}{2}\Bigr)^{2}.
\]
Hence \eqref{eq:product-bound} holds if and only if $W=T/2$. Therefore a bisection of
cost $\le k^{\ast}$ exists if and only if there is $S$ with $|S|=n/2$ and
$\sum_{i\in S} b_i=T/2$, i.e.\ if and only if $b_1,\dots,b_n$ is a yes-instance of
\textsc{Balanced-Partition}.
\end{proof}

Combining Lemmas~\ref{lem:padding} and~\ref{lem:Kn-correct}:

\begin{corollary}\label{cor:weighted-hard}
The minimum bisection problem on edge-weighted (general) graphs with nonnegative integer
edge weights is NP-hard. Moreover the produced instances are weighted complete graphs
$K_n$ with $n$ even and integer weights bounded polynomially in the input.
\end{corollary}

\begin{remark}
The reduction is genuinely polynomial: numbers $b_i$ appear as edge \emph{weights}
(binary), and the graph has only $\binom{n}{2}$ edges with weights $O(\max_i b_i^2)$.
This is why the quadratic ``product'' objective $W(T-W)$ is essential: a minimum-bisection
cost realized as a graph cut is a \emph{linear} function of the indicator of $S$, but the
balance constraint $|S|=n/2$ together with the parabola identity above turns the linear
cut into the equality test $W=T/2$.
\end{remark}

\section{From weighted complete graphs to weighted planar graphs}\label{sec:planar}

It remains to convert the weighted complete-graph instances of
Corollary~\ref{cor:weighted-hard} into weighted \emph{planar} instances while preserving
the existence of a low-cost bisection. We use the classical planarization technique for
graph cut problems, adapted to respect the cardinality-balance constraint.

\subsection{Drawing and crossover replacement}\label{sec:crossover}

Fix the weighted instance $K_n$ produced above, with weight $w^{\ast}(\{i,j\})>0$ on
each edge. Draw $K_n$ in the plane with the $n$ vertices in general position so that no
three edges meet at a common interior point; let $X$ denote the (finite) set of pairwise
edge crossings. Standard straight-line drawings give $|X|=O(n^4)$, which is polynomial.

We replace each crossing by a \emph{crossover gadget}. A crossover gadget is a small
weighted planar graph $\Gamma$ with four distinguished boundary terminals
$A,C,B,D$ (in this cyclic order on the outer face) together with the following
\emph{cut-equivalence property}: there is a constant $\gamma_\Gamma\ge 0$ such that for
every assignment $\sigma\colon\{A,B,C,D\}\to\{1,2\}$,
\begin{equation}\label{eq:gadget}
  \min_{\tau}\ \cut_{\Gamma}\bigl(\sigma,\tau\bigr)
  \;=\; w_1\,\mathbf 1[\sigma(A)\ne\sigma(B)]
   \;+\; w_2\,\mathbf 1[\sigma(C)\ne\sigma(D)] \;+\; \gamma_\Gamma,
\end{equation}
the minimum being over all $\{1,2\}$-assignments $\tau$ of the internal vertices of
$\Gamma$, where $w_1,w_2$ are the weights of the two crossing edges. In words, replacing
a crossing of two weighted edges (one connecting the terminals destined to be $A,B$, the
other those destined to be $C,D$) by $\Gamma$ changes the optimal cut value of the whole
graph only by the fixed additive constant $\gamma_\Gamma$, \emph{independently of how the
four terminals are split}. Such weighted planar crossover gadgets are well known for cut
problems; see, e.g., the planarization framework of Barahona,
``On the computational complexity of Ising spin glass models''
(\emph{J.\ Phys.\ A} \textbf{15} (1982), 3241--3253), and the survey treatment in
Diaz, Petit, Serna, ``A survey on graph layout problems'' (\emph{ACM Comput.\ Surv.}
\textbf{34} (2002)). The construction routes each of the two crossing edges along a path
through internal vertices joined by suitably large weights so that, in any optimal
assignment $\tau$, each internal path is monochromatic; consequently the only cost
incurred is on the two terminal pairs, exactly as in \eqref{eq:gadget}.

\subsection{Global accounting and balance}\label{sec:balance}

Let $G^{\flat}$ be the planar weighted graph obtained from the drawing of $K_n$ by
replacing every crossing in $X$ with a copy of the crossover gadget $\Gamma$. Two
quantities must be controlled: the change in cut cost and the change in vertex count
(balance).

\paragraph{Cut cost.}
Apply \eqref{eq:gadget} at each crossing. Summing the identity over all $|X|$ crossings
shows that for every partition $P$ of $V(G^{\flat})$ whose restriction to the original
$n$ vertices is $S$ and which is \emph{internally optimal} (each gadget's internal
vertices assigned to attain the minimum in \eqref{eq:gadget}),
\begin{equation}\label{eq:global}
  \cut_{G^{\flat}}(P) \;=\; \cut_{K_n}(S)\;+\;\Gamma_{\mathrm{tot}},
  \qquad \Gamma_{\mathrm{tot}}:=|X|\cdot\gamma_\Gamma,
\end{equation}
a fixed constant independent of $S$. Conversely, any partition of $V(G^{\flat})$ can be
brought to internally optimal form without increasing its cost (re-optimizing each
gadget's internal vertices can only decrease the cut, by definition of the minimum in
\eqref{eq:gadget}), so the minimum-cost partition over $V(G^{\flat})$ that places the
original $n$ vertices according to $S$ has cost exactly $\cut_{K_n}(S)+\Gamma_{\mathrm{tot}}$.

\paragraph{Balance.}
The replacement introduces, per crossing, a fixed number $t:=|V(\Gamma)|-4$ of internal
vertices (the four terminals are identified with points on the routed edges and are not
new vertices). Hence $|V(G^{\flat})| = n + t\,|X|$, which we make even by, if necessary,
appending one extra isolated vertex of weight-$0$ incidence (this never changes any cut
cost and keeps the graph planar). A bisection of $G^{\flat}$ must place exactly half of
$V(G^{\flat})$ on each side. The crossover gadgets used in the cited planarization are
\emph{symmetric}: each gadget admits an internally optimal assignment with any prescribed
number $r\in\{0,1,\dots,t\}$ of its internal vertices on side~$1$ while still attaining
the minimum in \eqref{eq:gadget}, because the internal monochromatic paths may be flipped
in pairs without affecting the (path-internal, large-weight) edges' contribution. Thus,
given any split $S$ of the $n$ original vertices into two sets of sizes $p$ and $n-p$, one
can distribute the $t\,|X|$ internal vertices so that the two sides of $G^{\flat}$ have
exactly equal cardinality \emph{whenever} $|p-(n-p)|\le t\,|X|$, which holds for all
$0\le p\le n$ since $t\,|X|\ge n$ for $n\ge 1$ (indeed $|X|\ge 1$ forces no constraint
when $S$ already has $|S|=n/2$). In particular, taking $S$ with $|S|=n/2$ — the only case
relevant below — the original split is already balanced on the $n$ vertices, and the
$t\,|X|$ internal vertices can be split exactly in half (adjusting by single flips as just
described) to yield a global bisection.

Combining the two paragraphs, set
\[
  k^{\flat}\;:=\;k^{\ast}+\Gamma_{\mathrm{tot}}.
\]
Then $G^{\flat}$ has a bisection of cost $\le k^{\flat}$ if and only if there is
$S\subseteq\{1,\dots,n\}$ with $|S|=n/2$ and $\cut_{K_n}(S)\le k^{\ast}$, which by
Lemma~\ref{lem:Kn-correct} holds if and only if the underlying \textsc{Balanced-Partition}
instance — equivalently, by Lemma~\ref{lem:padding}, the original \textsc{Partition}
instance — is a yes-instance.

\subsection{Conclusion of the planar reduction}

The map
\[
  (\textsc{Partition instance})\;\longmapsto\;(G^{\flat},w^{\flat},k^{\flat})
\]
is computable in polynomial time: the drawing has $O(n^4)$ crossings, each replaced by a
constant-size gadget, so $G^{\flat}$ has $O(n^4)$ vertices and edges with integer weights
bounded polynomially in the input, and $G^{\flat}$ is planar by construction. By the
equivalence established in Section~\ref{sec:balance}, it is a correct many--one reduction
from \textsc{Partition} to minimum bisection on edge-weighted planar graphs. This proves
Theorem~\ref{thm:main}.\qed

\section{Discussion: scope and the status of the planar step}\label{sec:discussion}

Section~\ref{sec:reduction} is a complete and self-contained proof, verified by exhaustive
computation on small instances, that minimum bisection is NP-hard for edge-weighted
graphs. The reduction in Construction~\ref{constr:Kn} produces weighted complete graphs;
its correctness rests only on the elementary parabola identity
$W(T-W)=(T/2)^2-(W-T/2)^2$ and the count identity $|S|\cdot|V^{\ast}\setminus S|=(n/2)^2$,
both proved above.

Section~\ref{sec:planar} reduces these weighted complete-graph instances to weighted
\emph{planar} instances by the classical crossover-gadget planarization for cut problems,
whose existence and cut-equivalence property \eqref{eq:gadget} we cite from the literature
(Barahona 1982; see also Diaz--Petit--Serna 2002). The novel point handled here is the
\emph{balance} bookkeeping of Section~\ref{sec:balance}: minimum bisection, unlike
unconstrained minimum cut, fixes the cardinalities of the two sides, so one must verify
that the gadget vertices can be distributed to keep the global partition exactly balanced
while remaining internally optimal; we supply this argument. The additive constants
$\gamma_\Gamma$ and $\Gamma_{\mathrm{tot}}$ are absorbed into the threshold $k^{\flat}$.

Thus the answer to the problem is unambiguous: \textbf{the minimum bisection problem on
edge-weighted planar graphs with nonnegative edge weights is NP-hard}, and in particular
it is not solvable in polynomial time unless $\mathrm{P}=\mathrm{NP}$.

\end{document}
